\documentclass[11pt]{article}
\input{preamble}
\externaldocument[][nocite]{revision_2026-10-16/build/supplement}[supplement.pdf]
% Set \anonymoustrue for an author-free review copy if required by the portal.
\newif\ifanonymous
\anonymousfalse
\title{The Cost of Reliable Near-Optimal Set Maintenance\\with Arriving Arms}
\ifanonymous
\author{}
\else
\author{Deqi Zheng\thanks{The first two authors contributed equally.} \quad
 Xiaoyang Xu\footnotemark[1] \\ Qiuzhen College, Tsinghua University
 \and Yuhong Yang\thanks{Corresponding author: yyangsc@mail.tsinghua.edu.cn.}\\
 Yau Mathematical Sciences Center, Tsinghua University}
\fi
\date{October 15, 2026}
\ifanonymous
\hypersetup{pdftitle={The Cost of Reliable Near-Optimal Set Maintenance with Arriving Arms},pdfauthor={}}
\else
\hypersetup{pdftitle={The Cost of Reliable Near-Optimal Set Maintenance with Arriving Arms},pdfauthor={Deqi Zheng, Xiaoyang Xu, Yuhong Yang}}
\fi
\begin{document}
\maketitle

\begin{abstract}
We study the cost of maintaining reliable reports of all relatively
near-optimal arms while collecting rewards when new arms arrive and remain
available. Reward laws are stationary, but each arrival can change the
membership boundary. We measure reporting progress by cumulative unresolved
arm time outside a prescribed boundary band, under simultaneous correctness
of all issued labels. Uncertainty about one new arm can obstruct positive
certification of many earlier candidates. A lower bound in one common
reference environment quantifies this effect. We obtain matching
worst-schedule orders for unresolved area and an arrival-growth boundary
for jointly sublinear area and expected regret. On a specified
repeated-update family, rapid maintenance requires inverse-square regret,
whereas a constant-factor larger area budget admits a logarithmic regret
upper bound. The rapid cost also holds with nonvanishing probability in a
fixed environment. A fixed-precision checkpoint construction provides the
upper bounds without knowing future arrivals. Finite-budget comparisons
show that its conservative certification cost can be substantial; they
do not establish uniform numerical superiority or a general optimal
regret--area frontier.
\end{abstract}
\noindent\textit{Keywords:} sequential experimentation; arriving arms;
near-optimal set maintenance; dynamic regret; simultaneous inference.

\section{Introduction}
\label{sec:introduction}

Sequential experiments often acquire new candidates while earlier ones remain
available. An allocation rule must learn which candidates yield high rewards.
An experimenter may also need to know which of all available candidates are
within a specified tolerance of the current best. These objectives need not
favor the same observations. A candidate can be sufficiently inferior to
receive little reward-driven exploration and still be difficult to classify
relative to the near-optimality boundary.

We study these objectives together under permanent exogenous arrivals and
stationary reward laws. If $r^*(t)$ is the largest mean among the arms
available at time $t$, an arm is relatively near-optimal when its mean is at
least $r^*(t)-\varepsilon$. Reports certify membership, certify exclusion, or
leave the arm unresolved. We require one simultaneous correctness event for
all labels at all times. To rule out indefinite abstention, we charge the
cumulative unresolved time of arms whose means lie more than $\eta$ from
the membership boundary. The question is how much regret is necessary, and
how much is attainable, under this reporting requirement.

New arrivals affect the task in two ways. An uncertain new candidate may be
better than the incumbent by enough to invalidate many existing positive
reports. Evidence about that one candidate is then relevant to many
reports at once. A later improvement can also make an earlier difficult
membership question unnecessary: a candidate close to the old boundary may
be far below the new one. Its existing upper confidence bound can be
combined with a new reference's lower bound to certify exclusion. Neither
effect requires the reward law of an old arm to change.

Our observation model and regret comparator have established precedents.
Permanent arrivals form a monotone special case of sleeping-bandit
availability \citep{kleinberg2010}; ballooning bandits already compare reward
with the best arm currently available \citep{ghalme2021}. Exogenous arrivals
differ from drawing arms on demand from a reservoir
\citep{10.1214/aos/1069362389} and from the disappearing actions of mortal
bandits \citep{NIPS2008_788d9869}. Benchmark improvements here arise from new
arms, whereas nonstationary bandits allow the reward distributions of
existing arms to change \citep{Besbes2014,Auer2019ADSWITCH}. Sequential
allocation and adaptive experimentation have a long statistical history
\citep{Robbins1952,LAI19854,Lai1987}; the expanding set of actions retains
that allocation problem and adds an evolving inferential target.

For a fixed arm set, identifying every relatively near-optimal arm is the
ALL-$\varepsilon$ problem \citep{mason2020}. Joint reward and identification
objectives also have established precedents
\citep{degenne2019,yang2026,cho2025}. Our task requires continuing reports
under arrivals, with all-time label reliability and an explicit cumulative
cost for unresolved arm time. The distinction concerns the reporting
protocol and its cost, rather than a new arriving-arm model. The simultaneous
recovery scales used in our lower bounds build on existing information
inequalities and multiple-recovery arguments
\citep{kaufmann2016,malloy2014,jamieson2018}.

We give three results about maintenance cost. First, a common-reference
lower bound shows how one uncertain arrival can obstruct many other
certifications. Its information queue accumulates cost over time within
one reward environment. Second, we obtain matching worst-schedule orders
for cumulative unresolved area, including finite-horizon truncation, and
derive the polynomial arrival boundary for jointly sublinear area and
expected regret. This worst-case lower bound already occurs in an initial
static batch; the common-reference result describes a separate consequence
of continuing arrivals. Third, on a specified repeated-update family,
rapid resolution of temporary membership tasks has matching inverse-square
regret costs. Allowing a constant-factor larger area budget admits a
logarithmic regret upper bound. A nonvanishing-probability lower bound
holds in one fixed environment under the same cumulative-area requirement.
The deferred endpoint has only an upper bound, and the general frontier
over intermediate budgets and overlapping arrivals remains open.

An explicit fixed-precision procedure, Checkpoint-UCB, supplies the
attainability results. It shares all observations between certification
and reward learning and retains every arrived arm in reward selection.
The construction separates the statistical costs cleanly, but its
conservative quotas can exceed a practical budget. Our numerical study
retains those failures as well as favorable comparisons. Section~\ref{sec:model}
defines the joint task, Section~\ref{sec:method} gives the construction,
and Sections~\ref{sec:maintenance}--\ref{sec:tradeoff} establish the
maintenance and regret bounds. Section~\ref{sec:experiments} examines their
finite-budget implications.

\section{Model, reliable reports, and maintenance cost}
\label{sec:model}
\label{sec:problem}

\subsection{Observation model and regret}

Let $\mathcal A$ be a possibly infinite set of arms. Arm $i$ arrives at
time $\tau_i$, after which it remains available. At every physical time
$t=1,2,\ldots$,
\[
 \mathcal A_t=\{i:\tau_i\le t\},\qquad
 1\le K_t=|\mathcal A_t|<\infty.
\]
Arbitrary finite batches are allowed, including $K_T>T$ at an evaluation
endpoint. Each arm has an unknown stationary reward law on $[0,1]$ with
mean $r_i$. Its successive potential observations $(Y_{i,n})_{n\ge1}$ are
independent and identically distributed, and the streams of different arms
are independent. Arrivals are exogenous and independent of both these
streams and policy randomization. Unless stated otherwise, each result
is for a fixed deterministic arrival schedule.

At time $t$, the policy first receives all current arrivals, then selects
one arm $I_t\in\mathcal A_t$, observes one reward $X_{t,I_t}$, and finally
issues its report. It uses the observed history
\[
 H_{t-1}=(\mathcal A_s,I_s,X_{s,I_s})_{s=1}^{t-1}
\]
and independent randomization. Means, future arrivals, and true membership
labels are unavailable to the policy. With
$N_i(t)=\sum_{s\le t}\ind\{I_s=i\}$, the physical observation is
$X_{t,I_t}=Y_{I_t,N_{I_t}(t-1)+1}$. Write
\[
 \hat r_i(n)=\frac1n\sum_{s=1}^nY_{i,s},\qquad
 r^*(t)=\max_{i\in\mathcal A_t}r_i,\qquad
 \Delta_i(t)=r^*(t)-r_i \quad(i\in\mathcal A_t).
\]
At a deterministic endpoint $T$, cumulative dynamic regret is
\[
 R_T(\pi)=\sum_{t=1}^{T}\Delta_{I_t}(t).
\]
This is a sum of mean gaps, random through the policy's arm selections;
$\E R_T$ is its expectation. The quantity does not compare realized
reward sequences. Physical time always counts actual observations.

\subsection{Reliable reports and cumulative unresolved time}
\label{sec:reporting_task}

For $0<\varepsilon<1$, the target set is
\[
 \mathcal G_t^\varepsilon=\{i\in\mathcal A_t:\Delta_i(t)\le\varepsilon\}.
\]
After the current observation, a report partitions all of $\mathcal A_t$
into $\mathcal P_t,\mathcal N_t,\mathcal V_t$: certified members, certified
nonmembers, and unresolved arms. Define
\[
 E_{\rm lab}=\{\text{for all }t:\ \mathcal P_t\subseteq\mathcal G_t^\varepsilon
                  \subseteq\mathcal A_t\setminus\mathcal N_t\}.
\]
A policy is $\delta$-reliable if $\Prob(E_{\rm lab})\ge1-\delta$ in every
allowed reward environment. A positive report may be withdrawn after a new
arrival. On the correctness event, a negative report remains valid because
$r^*(t)$ is nondecreasing and the arm's mean is stationary.

Reliability permits all arms to remain unresolved. For a boundary width
$0<\eta<1$, define
\begin{equation}
 \begin{split}
 d_i^\varepsilon(t)&=|\Delta_i(t)-\varepsilon|,\\
 U_t(\eta)&=\sum_{i\in\mathcal A_t}
   \ind\{i\in\mathcal V_t,d_i^\varepsilon(t)>\eta\},\qquad
 \mathcal U_T(\eta)=\sum_{t=1}^{T}U_t(\eta).
 \end{split}
 \label{eq:area}
\end{equation}
The count includes newly arrived arms with no observations and old arms
that become unresolved again. Every such arm--time pair is charged; the
loss is not limited to the time before a first decision. Within the
boundary band, unresolved time is exempt, but any label issued there must
still be correct. The true distance $d_i^\varepsilon(t)$ is used only by
the evaluator and in the analysis.

The task parameters are fixed when comparing policies. Their roles are
summarized below.
\begin{center}
\begin{tabular}{cl}
\toprule
Parameter & Meaning\\
\midrule
$\varepsilon\in(0,1)$ & Relative near-optimality tolerance\\
$\delta\in(0,1/2)$ & Joint error budget\\
$\eta\in(0,1)$ & Boundary width exempting unresolved time\\
$\rho\in(0,1)$ & Predetermined precision of reporting checkpoints\\
\bottomrule
\end{tabular}
\end{center}

The main requirement at a fixed $T$ is
\begin{equation}
 \E_\theta R_T\le r,\qquad
 \Prob_\theta\{E_{\rm lab},\mathcal U_T(\eta)\le a\}\ge1-\delta.
 \label{eq:joint_protocol}
\end{equation}
It combines an expected regret bound with a joint high-probability
label--area guarantee. In particular, $\mathcal U_T=o(T)$ requires the
time-average number of unresolved separated arms to vanish, rather than
only a small unresolved fraction of all arms. The area guarantee is for
each deterministic endpoint; it need not hold for every endpoint on one
common area event.

For a performance family $\Theta$, let $\mathfrak R_\Theta(a)$ be the
infimum of $\sup_{\theta\in\Theta}\E_\theta R_T$ over policies reliable on
the full bounded model and satisfying the probability condition in
\eqref{eq:joint_protocol} for every $\theta\in\Theta$. The value is
$+\infty$ if the feasible class is empty. Full-model label reliability
and a family-specific area guarantee are distinct requirements. The
lower bounds remain valid when reliability is required only on the
lower-bound subfamily.

\subsection{How arrivals change the reporting problem}

For a fixed old arm, $\Delta_i(t)$ is nondecreasing, so its true membership
can change only from member to nonmember. Its boundary distance
$|\Delta_i(t)-\varepsilon|$ need not be monotone. A new candidate can first
create uncertainty about the current best mean, obstructing positive
certification of many old arms even when those arms remain truly
near-optimal. A sufficiently good new reference can later move the
boundary past earlier candidates, making their fine classification
unnecessary. Sections~\ref{sec:maintenance} and~\ref{sec:tradeoff} quantify
these two effects under the same reporting protocol.
